% =====================================================================
%  4.7 Amenazas a la validez.  ESTADO: andamiaje.
%  Sección sensible: redactar tras cerrar §4.6 (resultados).
% =====================================================================
\section{Amenazas a la validez}
\label{cap:amenazas}

\pendienteredaccion{} Redactar tras consolidar los resultados. Estructura
prevista (por tipo de validez):

\subsection{Validez de constructo.}
Proxy de complejidad cognitiva $\neq$ SonarQube oficial; validación
cruzada limitada a un subset (\autoref{sec:sonar-val}). El oráculo de RQ3
se basa en la CC y en la comparación con el baseline determinista, no en
revisión humana ni en ejecución de los tests del proyecto original.

\subsection{Validez interna.}
Determinismo del transformador validado por tests; los LLMs presentan
no-determinismo residual aun con $T=0$ (mitigado con 3 intentos). La
política \emph{one-at-a-time} limita interacciones inesperadas pero podría
ocultar efectos compuestos.

\subsection{Validez externa.}
El corpus consolidado (126 casos, 104 elegibles) procede mayoritariamente
de proyectos abiertos, pero su composición no es una muestra
estadísticamente representativa del código Java «en general»: es una
\emph{muestra intencional} de métodos con condicionales anidados (sesgo de
selección). A ello se añade una limitación de reproducibilidad: parte de
los proyectos fuente no tiene versión ni \emph{commit} anclados y el corpus
mezcla versiones de un mismo proyecto (\autoref{sec:corpus}), por lo que
reconstruirlo idéntico exigiría reanclar dichas versiones. El subset de
RQ3 ($n=50$) cubre los ejes de variación del corpus, pero sigue siendo una
muestra intencional y no una muestra aleatoria representativa del código Java
en general. Además, la comparativa de RQ3 se restringe a dos modelos del mismo
proveedor (OpenAI):
\hipotesis{} su generalización a otras familias de modelos está por
comprobar. Por último, el patrón objetivo se limita al caso base sin
\lstinline|else|/\lstinline|else if|, por lo que los resultados no son
extrapolables a refactorizaciones más complejas.

\subsection{Validez de conclusión.}
Aunque el subset de RQ3 se ha ampliado a $n=50$ casos, sigue sin sostener
inferencia estadística fuerte sobre poblaciones de código; métrica única (CC)
que puede sesgar la valoración global de calidad.

\subsection{Mitigaciones implementadas.}
Congelación del protocolo (\texttt{LlmExperimentProtocol}); trazabilidad
completa (evidencias JSON/CSV/MD por invocación); guarda de elegibilidad
v1.1 en el oráculo. \pendientemetodo{} Decidir si se amplía a otro
proveedor de LLM o se documenta la restricción de forma definitiva.
